\scanpage{269}
\mspage{219}
\[
\mathrm{Ho}_W(S)\overset{\text{déf}}{=}(M_{/S})(W_S^0)^{-1}.
\]
Qu'il n'y ait pas de choix \uline{canonique} qui s'impose ici, déduit de la seule structure $(M,W)$, est clair sur l'exemple de $(\mathrm{Cat},W)$, auquel cas on trouve \uline{deux} choix symétriques transformés l'un en l'autre par l'involution
\[
X\longrightarrow X^{\unsure{\circ|0}}\qquad(\mathrm{Cat},W)\xrightarrow{\unsure{\sim}}(\mathrm{Cat},W),
\]
à savoir $W_S^g$ et $W_S^d$. Le choix de l'un ou de l'autre, disons (pour fixer les idées) de $W_S^g$, est associé à une notion d'\uline{iso local}\unsure{``|''} (ou de ``\unsure{coïso|coiso|coïsom.} local\unsure{''|.} dans le cas de $W_S^d$) que je voudrais axiomatiser. Il semble bien que pour les principales catégories de modèles pour $\mathrm{Hot}$ (et pour ses variantes), on dispose de façon assez naturelle d'une telle notion d'iso local (sans l'embarras, le plus souvent, d'un \unsure{autom.|autome} involutif fournissant une notion symétrique \ldots)

On supposera donc, comme dans loc. cit. p. 94-\unsure{95|96}, donné une classe
\[
\unsure{\mathcal L|L}\subset\mathrm{Fl}(M)
\]
de morphismes, jouant le rôle de ``flèches de
% archive sheet number: 269
